\section{Background: Three Governance Frameworks}
\label{sec:background}

Three governance instruments, despite differing in legal status and geography, converge on a common demand: that claims about a high-risk AI system's trustworthiness be supported by evidence. We introduce each in turn---the US National Institute of Standards and Technology (NIST) AI Risk Management Framework, the European Union's binding regulation, and the international management-system standard---then distill the shared lifecycle and risk spine that motivates our methods-to-requirements mapping in Table~\ref{tab:matrix}.

\subsection{The NIST AI Risk Management Framework}
\label{subsec:nist}
The NIST AI RMF~1.0 (January 2023) is a voluntary framework to help organizations manage risks across the AI lifecycle~\cite{nist2023airmf}. It is structured around four interacting functions. \textit{GOVERN} establishes the organizational culture, policies, and accountability structures that cut across the other functions; \textit{MAP} characterizes the context, intended purpose, technical methods, and potential impacts of a system; \textit{MEASURE} applies quantitative and qualitative techniques to analyze and track identified risks; and \textit{MANAGE} allocates resources and selects response strategies for prioritized risks~\cite{nist2023airmf}. The MEASURE function is the most directly technical: its MEASURE~2.x subcategories require that test sets, metrics, and tooling for assessment, audit, verification, and validation be documented (MEASURE~2.1), that evaluations involving human subjects meet applicable requirements and are representative of the relevant population (MEASURE~2.2), and that performance and assurance criteria be measured under deployment-like conditions (MEASURE~2.3); further MEASURE~2.x subcategories address safety (2.6), security and resilience (2.7), explainability (2.9), privacy (2.10), and fairness and harmful bias (2.11). These operationalize the framework's seven trustworthiness characteristics: validity and reliability, safety, security and resilience, accountability and transparency, explainability and interpretability, privacy, and managed fairness~\cite{nist2023airmf}. NIST extended this scaffolding with a Generative AI Profile that enumerates risks specific to generative systems, such as confabulation, dangerous content, and information integrity, and maps suggested actions back to the four functions~\cite{nist2024genai}. U.S. executive-branch AI policy has shifted across administrations; the AI RMF's voluntary status means its adoption is driven by organizational choice and sectoral expectation rather than statutory mandate.

\subsection{The EU AI Act High-Risk Regime}
\label{subsec:euaiact}
Regulation (EU) 2024/1689, the EU AI Act, is a binding, horizontally applicable law that classifies certain systems as high-risk and imposes obligations on their providers~\cite{euaiact2024}. The core technical obligations are concentrated in Chapter~III, Section~2. Article~9 requires a risk management system maintained throughout the lifecycle; Article~10 imposes data and data governance requirements on training, validation, and testing data sets; Article~11 mandates technical documentation containing the elements of Annex~IV; Article~12 requires record-keeping and automatic logging; Article~13 governs transparency and the provision of information to deployers; Article~14 requires human oversight through appropriate human-machine interface tools; and Article~15 requires appropriate levels of accuracy, robustness, and cybersecurity; its cybersecurity provision names resilience against attempts to alter use or behavior by exploiting system vulnerabilities, and---where appropriate---measures to prevent, detect, respond to, resolve, and control attacks such as data and model poisoning, adversarial examples (model evasion), confidentiality attacks (e.g., model inversion and extraction), and model flaws~\cite{euaiact2024}.

Compliance is not self-asserted in the abstract: Article~43 establishes conformity assessment procedures under which technical documentation, the quality management system of Article~17, and risk management arrangements are verified, in defined cases by a notified body, before market placement. Article~72 then requires providers to establish a post-market monitoring system that collects and analyzes performance and incident data across the operational life of the system~\cite{euaiact2024}. The interpretive burden of translating these legal terms into testable engineering requirements has motivated recent work on concept and requirement mapping between the Act and international standards~\cite{hernandez2024knowledgegraph,rintamaki2025ontology,lucaj2025techops}, on the standing of fundamental rights within harmonized standardization~\cite{hodac2024fundamentalrights}, and on the relationship between algorithmic fairness techniques and the Act's non-discrimination provisions~\cite{meding2025algorithmicdiscrimination}.

\paragraph{Application timeline (as of June 2026)} The obligations above are \emph{enacted} in
Regulation~(EU)~2024/1689, but their \emph{application dates} have been revised, and we are
careful to distinguish enacted text, currently applicable provisions, and agreed-but-not-yet-consolidated
amendments. Provisions already applicable include the prohibited-practices and AI-literacy duties
(from February~2025) and the general-purpose-AI model obligations (from August~2025). A political
agreement reached on 7~May~2026 (the European Commission's ``Digital Omnibus'') then deferred the
high-risk obligations: the Annex~III (use-case) high-risk requirements move from 2~August~2026 to
\textbf{2~December~2027}, and the Annex~I (product-embedded) high-risk requirements from
2~August~2027 to \textbf{2~August~2028}; a blanket ``stop-the-clock'' moratorium was rejected. As
these amendments were, at the time of writing, agreed but not yet consolidated into the cited
regulation, we treat them as anticipated rather than enacted law and avoid presenting the high-risk
regime as an immediate compliance deadline.

\subsection{ISO/IEC 42001:2023 and the AI Management System}
\label{subsec:iso}
ISO/IEC 42001:2023 specifies requirements for an Artificial Intelligence Management System (AIMS) and follows the harmonized high-level structure shared across ISO management-system standards~\cite{iso42001}. Clauses~4 through~10 define the management-system requirements: context of the organization (Clause~4), leadership (Clause~5), planning (Clause~6), support (Clause~7), operation (Clause~8), performance evaluation (Clause~9), and improvement (Clause~10). The normative Annex~A provides a catalogue of controls (with the corresponding implementation guidance given in the informative Annex~B), spanning AI policies (A.2), internal organization (A.3), resources for AI systems (A.4), assessing impacts of AI systems (A.5), the AI system lifecycle covering planning, design, development, testing, validation, deployment, operation, monitoring, maintenance, and decommissioning (A.6), data for AI systems (A.7), information for interested parties (A.8), responsible use of AI systems (A.9), and third-party and customer relationships (A.10)~\cite{iso42001}. Several Annex~A controls call for technical evidence rather than policy alone, notably A.6.2.4 on AI system verification and validation and A.6.2.6 on AI system operation and monitoring. The standard is designed to be used alongside ISO/IEC 23894:2023, which provides organizational guidance on AI-specific risk management and informs the planning and impact-assessment activities of the AIMS~\cite{iso23894}.

\subsection{A Shared Lifecycle and Risk Spine}
\label{subsec:spine}
Although the three instruments differ in legal force, scope, and vocabulary, they share a recurring structure: a risk-based lifecycle in which risks are identified and mapped, measured and evaluated, managed and mitigated, and monitored after deployment. The NIST MAP/MEASURE/MANAGE progression~\cite{nist2023airmf}, the EU AI Act's sequence from risk management (Article~9) through accuracy and robustness (Article~15) to post-market monitoring (Article~72)~\cite{euaiact2024}, and the ISO/IEC 42001 operation, performance-evaluation, and improvement clauses together with the A.6 lifecycle controls~\cite{iso42001} are mutually recognizable expressions of this spine. Critically, each framework reaches a point at which a governance claim can be discharged only through technical evaluation evidence: documented test sets and metrics for trustworthiness characteristics under NIST MEASURE~2.x~\cite{nist2023airmf}, the Annex~IV technical documentation and accuracy, robustness, and cybersecurity demonstrations required for conformity assessment under the EU AI Act~\cite{euaiact2024}, and the testing, validation, and monitoring records expected of an AIMS~\cite{iso42001}. This convergence on evaluation evidence, observed across comparative studies of regional frameworks~\cite{almaamari2025crossregional,gadekallu2025responsibleaisurvey} and assurance proposals~\cite{lupo2026trustworthyposture,puhlfurss2025modelcards}, is the gap this review addresses by mapping concrete technical assessment methods to the specific requirements they can satisfy. The shared spine should not be mistaken for equivalence, however: the three instruments differ materially in legal force, unit of assessment, responsible actor, and assurance mechanism (Table~\ref{tab:frameworks}). A single technical artifact may be \emph{relevant} to all three, but relevance, evidentiary sufficiency, and legal acceptability are distinct; in particular, ISO/IEC~42001 certification of an organization's management system does not establish that an individual model is EU~AI~Act compliant, and the NIST AI RMF is voluntary. We therefore frame our crosswalk as a map of \emph{technical relevance}, not of legal conformity.

% Framework non-equivalence table (reviewer R2).
\begin{table*}[t]
\centering
\caption{The three instruments are \emph{not} equivalent assessment objects. A shared
technical artifact may be relevant to all three, but their legal status, unit of assessment,
responsible actor, and assurance mechanism differ. In particular, ISO/IEC~42001 certification
certifies an organization's management system, not that an individual AI system is EU~AI~Act
compliant, and the NIST AI RMF is voluntary guidance, not law. \textbf{Takeaway:} the three
regimes overlap heavily in \emph{evidence} but differ in legal force, unit of assessment,
responsible actor, and assurance mechanism; shared evidence must not be read as legal
equivalence.}
\label{tab:frameworks}
\footnotesize
\setlength{\tabcolsep}{4pt}
\begin{tabular}{p{0.16\textwidth}p{0.25\textwidth}p{0.27\textwidth}p{0.24\textwidth}}
\toprule
Dimension & NIST AI RMF 1.0 & EU AI Act (Reg.\ (EU) 2024/1689) & ISO/IEC 42001:2023\\
\midrule
Legal status & Voluntary guidance & Binding law (phased application) & Voluntary, certifiable standard\\
Object assessed & AI risk-management \emph{practice} of an organization & The \emph{high-risk AI system}, per its intended purpose & The organization's \emph{AI management system} (AIMS)\\
Responsible actor & Any AI actor (adopter) & Provider (with deployer duties) & The organization seeking certification\\
Assurance mechanism & Self-directed; no external conformity assessment & Conformity assessment, self- or notified-body (Art.~43) & Third-party certification audit\\
Lifecycle role & GOVERN/MAP/MEASURE/MANAGE, lifecycle-wide & Pre-market \emph{and} post-market obligations (Arts.~9--15, 72) & Plan--Do--Check--Act management cycle\\
Role of technical evidence & Populates the MEASURE function & One input among QMS, documentation, and the CA route & Evidence that AIMS controls operate as intended\\
\bottomrule
\end{tabular}
\end{table*}

